\documentclass[12pt]{extarticle}
\def\activityid{F04-U-v2}\usepackage[letterpaper,margin=.65in,top=.60in,bottom=.65in]{geometry}
\usepackage[T1]{fontenc}\usepackage[default]{sourcesanspro}
\usepackage{microtype,tikz,array,tabularx,fancyhdr,amsmath,amssymb}
\usepackage[hidelinks]{hyperref}\usetikzlibrary{calc,arrows.meta}
\setlength{\parindent}{0pt}\setlength{\parskip}{7pt}\setlength{\headheight}{0pt}\setlength{\footskip}{26pt}
\setlength{\emergencystretch}{2em}\pagestyle{fancy}\fancyhf{}\renewcommand{\headrulewidth}{0pt}\renewcommand{\footrulewidth}{.4pt}
\fancyfoot[L]{\scriptsize Bellingham Math Circle / Week 4 / \activityid}\fancyfoot[R]{\scriptsize \thepage}
\definecolor{ink}{gray}{.12}\definecolor{soft}{gray}{.95}\color{ink}
\newcommand{\PageTitle}[2]{{\fontsize{10}{12}\selectfont\bfseries WEEK 4\quad ROUND AND ROUND\hfill #1\par}\vspace{3pt}{\fontsize{26}{29}\selectfont\bfseries #2\par}\vspace{2pt}}
\newcommand{\NameLine}{{\small Name: \rule{2.65in}{.4pt}\hfill Date: \rule{1.35in}{.4pt}\par}\vspace{4pt}}
\newcommand{\Task}[2]{\par\vspace{3pt}{\large\bfseries #1}\par #2\par}
\newcommand{\Subhead}[1]{\par\vspace{3pt}{\large\bfseries #1}\par}
\newcommand{\RuleBox}[1]{\par\begingroup\setlength{\fboxsep}{8pt}\colorbox{soft}{\parbox{\dimexpr\linewidth-16pt\relax}{#1}}\endgroup\par}
\newcommand{\WriteBox}[2]{\par\textbf{#1}\par\vspace{2pt}\begin{tikzpicture}\draw[black!40,line width=.5pt] (0,0) rectangle (\dimexpr\linewidth-.5pt\relax,#2);\end{tikzpicture}\par}
\newcommand{\StudentFooter}{\vfill{\small Try it with objects. Explain by pointing, talking, or drawing.}\par}
\newcommand{\LampRing}[3]{\begin{tikzpicture}[x=1in,y=1in]
\foreach \i in {1,...,#1}{\coordinate (v\i) at ({90-360*(\i-1)/#1}:#2);}
\foreach \i in {1,...,#1}{\pgfmathtruncatemacro{\j}{mod(\i,#1)+1}\draw[thick] (v\i)--node[midway,fill=white,inner sep=2pt,font=\small]{\i\j}(v\j);}
\foreach \i in {1,...,#1}{\node[circle,draw,fill=white,minimum size=.52in,inner sep=0pt,font=\bfseries] at(v\i){\i};}
\foreach \i in {#3}{\node[circle,draw,fill=ink,text=white,minimum size=.52in,inner sep=0pt,font=\bfseries] at(v\i){\i};}
\end{tikzpicture}}
\newcommand{\Lights}[1]{\begin{tikzpicture}[x=.55in,y=.55in,baseline=-.10in]\foreach \v [count=\i] in {#1}{\ifnum\v=1\fill (\i,0) circle(.16in);\else\draw (\i,0) circle(.16in);\fi}\end{tikzpicture}}
\newcommand{\ClockRing}[2]{\begin{tikzpicture}[x=1in,y=1in]\draw[black!30] (0,0) circle(#2);\pgfmathtruncatemacro{\n}{#1-1}\foreach\i in {0,...,\n}{\node[circle,draw,fill=white,minimum size=.35in,inner sep=0pt] at({90-360*\i/#1}:#2){\i};}\draw[-{Stealth},thick] (-.35,.15) arc(155:25:.38);\node[font=\small] at(0,-.18){clockwise};\end{tikzpicture}}
\newcommand{\Key}[2]{\begin{center}\renewcommand{\arraystretch}{1.55}\begin{tabular}{c|*{#1}{c}}in&#2\end{tabular}\end{center}}

\begin{document}

\PageTitle{GRADES 4--5 / 1}{The twelve-place code}\NameLine
\RuleBox{Use positions 0 through 11 clockwise. A shift +k moves every input k places clockwise, wrapping after 11. Decoding reverses that shift. A turn repeats the \textbf{same} shift.}
\begin{center}\ClockRing{12}{1.05}\end{center}
\Task{1. A broken message?}{Input 9 became output 1. Find the shift +k, where k is between 0 and 11. Check whether input 2 could become 6 under the same shift. Decode the output sequence \textbf{1, 4, 10}.}
\WriteBox{My key, check, and decoded message:}{.75in}
\Task{2. Follow one input.}{Starting at 0, repeat that shift until the first return. Then try starting at 1. Find \textbf{all} its separate loops, including every position exactly once.}
\WriteBox{All loops, and the first return time:}{.85in}
\StudentFooter
\newpage
\PageTitle{GRADES 4--5 / 2}{Classify every shift}
\Task{3. Find a rule, not just a table.}{Try enough shifts to spot a rule: start at 0, find the first positive return, and count separate loops. Use clockwise/counterclockwise symmetry. The full table is optional; move to page 3 when ready to explain.}
\begin{center}\renewcommand{\arraystretch}{1.45}\begin{tabular}{c|p{1.45in}|p{1.2in}|p{1.65in}}Shift&Turns to return&Number of loops&Visits all 12?\\\hline 0&&&\\1&&&\\2&&&\\3&&&\\4&&&\\5&&&\\6&&&\\7&&&\\8&&&\\9&&&\\10&&&\\11&&&\end{tabular}\end{center}
\Task{4. Predict before you test.}{On an 18-place ring, predict the return times for +6 and +5. What number connects the ring size, the step size, and the number of loops?}
\WriteBox{My conjecture and evidence:}{1.0in}
\StudentFooter
\newpage
\PageTitle{GRADES 4--5 / 3}{Why the rule works}
\Task{5. Count distance around the ring.}{After t turns of +k on an n-place ring, you have moved kt spaces. You are back at the start exactly when kt is a whole multiple of n. Use this to explain +4 on a 12-ring.}
\WriteBox{The first return and why none earlier works:}{.80in}
\Task{6. Shrink the common factor.}{For +8 on 12 places, return means 12 divides 8t, or 3 divides 2t. Draw three equal t-strips: together they total 3t. If two strips total a multiple of 3, why does removing them leave t a multiple of 3? Explain the first return.}
\WriteBox{My strip argument; optional general rule:}{.80in}
\Task{7. Undoing is different from visiting.}{The +4 rule does not visit every place from 0. Does it still give a reversible code on \emph{all} twelve inputs? Find its inverse. Then compare +4 followed by +7 with +7 followed by +4. Explain for every starting place.}
\WriteBox{Inverse, combined shift, and why the order agrees:}{.80in}
\textbf{Optional, with an adult:} Use the guide's subtraction-of-strips argument to justify coprime cancellation, then prove the gcd rule. Until then, label the general rule a conjecture or a supplied theorem.
\StudentFooter

\end{document}
